\documentclass[10pt,a4paper]{amsart}
\usepackage[margin=19mm]{geometry}
\usepackage{amsmath,amssymb,mathtools}
\usepackage[scr=rsfs,cal=euler]{mathalfa}
\usepackage{xfrac}
\usepackage[unicode,hidelinks,bookmarks=false]{hyperref}
\newcommand{\ie}{i.e.\ }
\newcommand{\cf}{cf.\ }
\newcommand{\resp}{respectively\ }
\newcommand{\Subsection}{\subsection}
\providecommand{\nobreakdash}{\nobreak\hbox{-}}
\setlength{\emergencystretch}{2em}
\multlinegap0pt
\usepackage[shortlabels]{enumitem}
\usepackage{amssymb}
\usepackage{xcolor}
\newtheorem{prop}{Proposition}
\title{$N$-Koszul algebras of finite global dimension for $N\geq 3$}
\author{So Nakamura}
\date{Source: arXiv:2608.20567v1 (CC BY 4.0)}
\begin{document}
\maketitle

\noindent\emph{Source and licence.} $N$-Koszul algebras of finite global dimension for $N\geq 3$. Version arXiv:2608.20567v1 (CC BY 4.0), distributed by arXiv under the Creative Commons Attribution 4.0 International licence, http://creativecommons.org/licenses/by/4.0/, as recorded in the arXiv licence metadata for this exact version and shown on the abstract page https://arxiv.org/abs/2608.20567v1. Full licence notice: \texttt{licenses/arxiv-2608.20567-CC-BY-4.0.txt}.\par\smallskip

\noindent\textit{Editorial scope.} Counted unit: the article's prop seqp (source0.tex lines 645--662). The article's complete original proof of it is delivered with it (source0.tex lines 664--762, one \texttt{proof} environment of 99 lines). No mathematics is written, repaired or re-derived: the statement and the proof are the article's own lines, in the article's own notation, and no retained line is substituted or re-flowed.  No further source-local context is needed: every cross-reference of every retained line resolves inside the two delivered spans.  Nothing was omitted from any retained span, so the package carries no omission record.  No retained line contains a citation, so the wrapper carries no bibliography and no bibliography entry.  No retained line contains a figure environment, an image-inclusion command or drawing code, so no image is delivered and the package declares no asset dependency.  Editorial formatting only, in addition: an AMS \texttt{amsart} preamble that restates the article's own declarations and macros used by the retained lines, namely the environments \texttt{prop} on separate counters, together with the standard packages those lines need.  The retained environments print in this document as Proposition 1 in order of appearance, whereas the article prints the same items with the numbering of its own full document.  The complete native source of the article as distributed in the arXiv e-print for this exact version is bundled unchanged as \texttt{sources/208186/source0.tex}.
\begin{prop}\label{seqp}
Let $s:=n_1=\pi_1$.
\begin{enumerate}
    \item We have $\pi_n=s^n$ for all $1\leq n\leq N-1$.
    \item $\alpha_1=\frac{s^N-s}{N}$. In particular, $s>1$.
    \item $\pi_{N+l}=\frac{-ls^{N+l}+(N+l)s^{l+1}}{N}+(N+l)s^{l-1}\alpha_{k-1}$ for all $1\leq l\leq N-1$.
    \item $\pi_{N+l}\geq\frac{N-(l+1)}{2N-1}s^{N+l}$ for all $1\leq l\leq N-1$.
    \item Let $N>3$. For all $N+2<n\leq 2N-1$, $\pi_n$ is divisible by $s^2$.
    \item Let $N>3$. For all $2N+3<n\leq 3N-1$, $\pi_n$ is divisible by $s^2$.

    
    %\item $m_{N+1}=\frac{-\alpha_k^{N+1}+(N+1)\alpha_k^2}{N}+(N+1)\alpha_{k-1}$.
    %\item $m_{N+2}=\frac{-2\alpha_k^{N+2}+(N+2)\alpha_k^3}{N}+(N+2)\alpha_k\alpha_{k-1}$.
    
    %\item $m_{N+1}\geq \frac{(p+1)(p-1)}{p(2p-1)}s^{p+2}-\frac{1}{p}s^{p+1}-\frac{p+1}{2p-1}s^3+\frac{p+1}{p}s^2$.
    %\item $m_{N+2}\geq \frac{p-3}{2p-1}s^{p+2}+\frac{(p+2)(p-1)}{p(2p-1)}s^3$.
\end{enumerate}
\end{prop}

\begin{proof}
$(1)$: Since $n<N$, $j_1=n$ and the other $j_l$'s are all $0$. Therefore $b_{n, i}=0$ if $i<N$. Hence  $\pi_n=(-1)^nb_{n, n}=(-1)^n(-\alpha_k)^n=s^n$.\\
$(2)$: If $n=N$, then $j_1=N$ or $j_N=1$ so that $b_{n, i}=0$ unless $i=1$ or $i=N$:
\begin{center}
\begin{tabular}{ |c|c|c|c|c|c|c|c| }
\hline
    $j_{N}$ & $j_1$ & $i$ & mult. coeff.\\
\hline
   $1$ &  & $1$ & $1$\\ 
\hline
      & $N$ & $N$ & $1$\\ 
\hline
\end{tabular}
\end{center}  
Therefore 
$$\pi_N=N\sum^N_{i=1}\frac{(-1)^i}{i}
b_{N, i}=N\left(-b_{N, 1}-\frac{b_{N, N}}{N}\right)=-N\alpha_1+\alpha_k^N=-N\alpha_1+s^N.$$
On the other hand, $\pi_N=n_1+Nn_N=n_1=s$ so that 
$$\alpha_1=\frac{s^N-s}{N}.$$
$s>1$ follows from the assumption that $\alpha_1>0$.\\

$(3)$: The following is the table for $n=N+l$ for every $1\leq l\leq N-1$:
\begin{center}
\begin{tabular}{ |c|c|c|c|c|c|c|c| }
\hline
  $j_{N+1}$ & $j_{N}$ & $j_1$ & $i$ & mult. coeff.\\
\hline
 $1$ & & $l-1$ & $l$ & $l$\\ 
\hline
    & $1$  & $l$ & $l+1$ & $l+1$\\ 
\hline
   &  & $N+l$ & $N+l$ & $1$\\ 
\hline
\end{tabular}
\end{center}  

Therefore 
\begin{eqnarray*}
    \pi_{N+l}&=&(N+l)\sum^{N+l}_{i=1}\frac{(-1)^i}{i}
b_{N+l, i}=(N+l)\left((-1)^lb_{N+l, l}+(-1)^{l+1}b_{N+l, l+1}+(-1)^{N+l}\frac{b_{N+l, N+l}}{N+l}\right)\\
    &=&(N+l)\left((-1)^l(-\alpha_k)^{l-1}(-\alpha_{k-1})+(-1)^{l+1}(-\alpha_k)^l\alpha_1+(-1)^{N+l}\frac{(-\alpha_k)^{N+l}}{N+l}\right)\\
    &=&(N+l)\alpha_k^{l-1}\alpha_{k-1}-(N+l)\alpha_k^l\alpha_1+\alpha_k^{N+l}\\
    &=&\frac{-(N+l)s^l(s^N-s)+Ns^{N+l}}{N}+(N+l)s^{l-1}\alpha_{k-1}\\
    &=&\frac{-ls^{N+l}+(N+l)s^{l+1}}{N}+(N+l)s^{l-1}\alpha_{k-1}\\
\end{eqnarray*}

$(4)$: Apply $l=N-1$ to $(3)$. Since $\pi_{2N-1}\geq0$, we obtain
$$(2N-1)s^{N-2}\alpha_{k-1}\geq\frac{(N-1)s^{2N-1}-(2N-1)s^N}{N}.$$
Therefore
$$\alpha_{k-1}\geq \frac{(N-1)s^{N+1}}{N(2N-1)}-\frac{s^2}{N}$$
therefore
\begin{eqnarray*}
m_{N+l}&=&\frac{-ls^{N+l}+(N+l)s^{l+1}}{N}+(N+l)s^{l-1}\alpha_{k-1}\\
&\geq& \frac{-ls^{N+l}+(N+l)s^{l+1}}{N}+(N+l)s^{l-1}\left(\frac{(N-1)s^{N+1}}{N(2N-1)}-\frac{s^2}{N}\right)\\
&=&\frac{(N+l)(N-1)-l(2N-1)}{N(2N-1)}s^{N+l}=\frac{N^2-(l+1)N}{N(2N-1)}s^{N+l}\\
&=&\frac{N-(l+1)}{2N-1}s^{N+l}
\end{eqnarray*}

$(5)$: Recall from $(3)$ that we have the following table for $n=N+l$ for every $1\leq l\leq N-1$:
\begin{center}
\begin{tabular}{ |c|c|c|c|c|c|c|c| }
\hline
  $j_{N+1}$ & $j_{N}$ & $j_1$ & $i$ & mult. coeff.\\
\hline
 $1$ & & $l-1$ & $l$ & $l$\\ 
\hline
    & $1$  & $l$ & $l+1$ & $l+1$\\ 
\hline
   &  & $N+l$ & $N+l$ & $1$\\ 
\hline
\end{tabular}
\end{center}  

It is clear that the coefficient of $\alpha_j$'s in each term of $\pi_{2N+l}$ viewed as a polynomial in $\alpha_j$'s is an integer. If $l>2$, then $j_1\geq2$ for each term of $\pi_n$. Thus $\pi_n$ is divisible by $\alpha_k^2=s^2$ if $2<l\leq N-1$.

$(6)$: The following is the table for $n=2N+l$ for some $2\leq l\leq N-1$:
\begin{center}
\begin{tabular}{ |c|c|c|c|c|c|c|c| }
\hline
 $j_{2N}$ & $j_{N+1}$ & $j_{N}$ & $j_1$ & $k$ & mult. coeff.\\
\hline
$1$ & & & $l$ & $l+1$ & $l+1$\\ 
\hline
 & $2$ & & $l-2$ & $l$ & $\binom{l}{2}$\\ 
\hline
  & $1$ & $1$ & $l-1$ & $l+1$ & $l(l+1)$\\ 
\hline
  & $1$ & & $N+l-1$ & $N+l$ & $N+l$\\ 
\hline
  & & $2$ &  $l$ & $l+2$ & $\binom{l+2}{2}$\\ 
\hline
  & & $1$ & $N+l$ & $N+l+1$ & $N+l+1$\\ 
\hline
  & &  & $2N+l$ & $2N+l$ & $1$\\ 
\hline
\end{tabular}
\end{center}  

It is clear that the coefficients of $\alpha_j$'s in the terms of $\pi_{2N+l}$ corresponding to $k=l+1, N+l, N+l+1 2N+l$ are all integers. If $k=l$, then the coefficient is $(2N+l)(l-1)/2$, and this is an integer since one of the two factors of the numerator is always even. If $k=l+2$, the coefficient is $(2N+l)(l+1)/2$ and this is also an integer. Therefore, $\pi_{n}$ can be viewed as a polynomial over $\alpha_j$'s with integer coefficients. If $l>3$, then $j_1\geq 2$ for each term of $\pi_n$. This means that every term of $\pi_n$ is divisible by $\alpha_k^2=s^2$ if $3<l\leq N-1$.  \end{proof}

\end{document}
